\section{build-pdfs and find-problems}\label{sec:ref-scripts}

Both scripts are documented in workflow terms in
section~\ref{sec:build-and-find}; this is the terse recap. Each prints its
full header with \texttt{-{}-help}.

\subsection*{build-pdfs}

\begin{manexsource}
build-pdfs [FILE.tex ...]    # no args: every .tex in the cwd
build-pdfs --release         # send due copies, build nothing
build-pdfs --release --dry-run   # ...and report instead of sending
build-pdfs --no-release FILE.tex # build without the end-of-run sweep
\end{manexsource}

\begin{center}
\begin{tabularx}{\linewidth}{@{}l Y@{}}
views & \texttt{\%!\ views: solutions, hints, projector, instructor} in
  the leading comment block; each produces \file{FILE-VIEW.pdf} by
  injecting its layer definition(s). Unknown views warn and skip. Says
  which PDFs a document has, never where they go.\\
copies & \texttt{\%!\ copies: default-->TAs, instructor-->TA-dir(from
  2026-10-01)} in the same block drops those PDFs, flat, into the named
  \file{.MIRRORDIR} destinations, on top of the mirror. \env{default} is
  the base PDF; any other name must also be on the views line.
  \texttt{(from \emph{date})} holds a copy until that day (inclusive):
  the build will not send it early, and the end-of-run sweep sends it
  once the date arrives. Anything malformed warns and sends nothing. A
  legacy \texttt{:LABEL} suffix on the views line warns and does not
  copy.\\
coverpage & \texttt{\%!\ coverpage: yes} in the same block also writes
  \file{FILE_COVERPAGE.pdf}, page~1 of the finished PDF on its own. From
  the base build only, never a view; same mirror handling as that PDF.
  \texttt{no} or no line is off, any other value warns. Needs
  \texttt{gs}; without it the extract is skipped and the build still
  succeeds.\\
output & nearest ancestor \file{.MIRRORDIR} (bare base path $+$ optional
  \texttt{BOTH}) tree-mirrors the repo under the base $\to$ else local.
  Base must exist; missing $\to$ warn and fall back to local.
  \texttt{LABEL: path} lines are the address book the \texttt{copies} row
  draws on (each must exist). \file{SAMPLE-MIRRORDIR} is the template.\\
aux files & \file{~/.cache/latexmk/}$\langle$\emph{repo-relative
  key}$\rangle$-$\langle$\emph{job}$\rangle$.\\
\file{.latexmkrc} & installed in each source's directory from
  \file{latexmkrc-shared} before compiling, when missing or out of date.
  A copy without the \texttt{managed by build-pdfs} marker line counts as
  hand-tuned and is left alone.\\
\end{tabularx}
\end{center}

\subsection*{build-images}

\begin{manexsource}
build-images [OPTIONS] FILE.tex ...
\end{manexsource}

One problem file in, one picture out --- for a learning-management system whose
question pools take an image rather than text. Each input yields
\file{NAME.pdf} and \file{NAME.png} (the question alone), plus
\file{NAME-SOLUTIONS.pdf} and \file{NAME-SOLUTIONS.png} (the question with its
worked solution), reusing build-pdfs' view suffix.

Those four are two independent choices, each defaulting to both, so any single
output can be asked for on its own: \texttt{-{}-no-question} and
\texttt{-{}-no-solutions} pick the variant, \texttt{-{}-no-pdf} and
\texttt{-{}-no-png} the format. \texttt{-{}-no-pdf} skips no work --- the PDF is
still built, in \env{\$TMPDIR}, because the PNG is rasterized from it --- it just
never reaches \texttt{-{}-out}. Turning off both variants, or both formats, is an
error rather than a silent no-op.

\begin{center}
\begin{tabularx}{\linewidth}{@{}l Y@{}}
wrapper & a throwaway \texttt{article} document in \env{\$TMPDIR}, never
  written into your tree. It loads \file{Exercises.sty} (and \file{tabularray},
  which every class carries, so a problem with a \env{tblr} builds) directly, because none
  of the three classes wants to be a single problem on a page: \file{Assessment}
  hangs numbers and marks in margin rails that fall outside a crop box, and
  \file{CourseDocument} brings a title block and headers a question has no use
  for.\\
heading & \cs{ProblemHeadingFormat} is emptied, so no ``Problem~N.'' appears
  --- the LMS numbers its own questions.\\
solutions & \cs{ShowSolutions} is defined \emph{before} \cs{documentclass}, the
  way build-pdfs injects it. Defined any later it is silently ignored
  (\file{Exercises.sty} reads it with \cs{ifdefined} at load time), and the
  solution is quietly absent.\\
\file{Workspace} & deliberately \emph{not} loaded, which is what makes the
  solution box render at its natural height: \file{Exercises} routes it through
  \cs{WSZeroHeight} only when that exists. Loading Workspace and calling
  \cs{collapseworkspace} reaches the same place by a longer road.\\
measure & \texttt{2.6\cs{alphabet}} --- the rule the classes use, not a number
  chosen here, so an image matches the same problem in a printed build and keeps
  matching if the typeface changes. \texttt{-{}-width} overrides it; a
  \env{matching} problem is the case that wants a wider one.\\
crop & vertical only. Every image comes out exactly one measure wide, so pool
  members display consistently wherever the LMS scales images to its column. The
  page is 200\,cm tall so that a long question can never split across two pages,
  which would yield two images of half a question each; the crop replaces the
  paper, so the height never reaches the output.\\
needs & \texttt{pdflatex}, \texttt{pdfcrop}, \texttt{gs}; \texttt{pdftoppm} for
  the PNG (\texttt{-{}-no-png} to skip).\\
\end{tabularx}
\end{center}

\subsection*{find-problems}

\begin{manexsource}
find-problems [--topic PAT] [--type PAT] [--difficulty PAT]
              [--source PAT] [--untagged] [--vocab] [--tsv]
              [--import [--with LIST] [--from DIR]] [PATH ...]
\end{manexsource}

\begin{center}
\begin{tabularx}{\linewidth}{@{}l Y@{}}
\texttt{-{}-topic}, \texttt{-{}-type}, \texttt{-{}-difficulty},
  \texttt{-{}-source} & AND-combined filters; PAT is a case-insensitive
  ERE matched per comma-separated item.\\
\texttt{-{}-vocab} & print the type/topic vocabulary with counts (filters
  ignored).\\
\texttt{-{}-untagged} & files with no \cs{ProblemMeta} block -- the
  backlog.\\
\texttt{-{}-import} & matches as paste-ready \cs{importproblem} lines.\\
\texttt{-{}-with LIST} & (with \texttt{-{}-import}) inject metadata keys as
  options, added where a file has them; \env{LIST} is a comma subset of the
  keys injectable for that bank -- \env{number,name} for problems,
  \env{kind,name,number} for facts. \env{points} stays manual.\\
\texttt{-{}-from DIR} & (with \texttt{-{}-import}) make each line's directory
  argument relative to \env{DIR}, the importing document's folder.\\
\texttt{-{}-tsv} & one tab-separated row per match
  (path, name, topics, type, difficulty, source, number) -- the machine feed
  for \file{pick-problems}. Respects the filters.\\
\end{tabularx}
\end{center}

A \env{PATH} operand may be a directory (walked recursively) or a single
\file{.tex} file, so a picker can ask \texttt{-{}-import} about exactly the
files it selected.

\subsection*{pick-problems}

\begin{manexsource}
pick-problems [--from DIR] [FILTER ...] [PATH ...]
\end{manexsource}

An fzf front-end to \file{find-problems}: it feeds \texttt{-{}-tsv} into fzf
(multi-select, metadata shown, the \file{.tex} previewed with \file{bat}), then
turns the chosen problems into paste-ready \cs{importproblem} lines via
\texttt{-{}-import}, printed to standard output. \texttt{-{}-from} is forwarded
to the emitter; every other argument is forwarded to the \texttt{-{}-tsv} feed,
so \file{find-problems}' filters pre-narrow what the picker offers. The Vim
command \texttt{:ImportProblem} (\texttt{<Space>ip}, in
\file{after/ftplugin/tex.vim}) is the same pipeline wired to insert at the
cursor, with \texttt{-{}-from} set to the editing buffer's own directory.

\subsection*{find-facts and pick-facts (the fact-bank twins)}

The fact bank (section~\ref{sec:ref-factmeta}) is searched and imported by the
same machinery, one rung smaller. \file{find-facts} and \file{pick-facts} behave
exactly like their problem counterparts above but read \cs{FactMeta} blocks and
emit \cs{importfact} lines, over the fact schema: the filters are
\texttt{-{}-topic} and \texttt{-{}-source} only (a fact has no type or
difficulty), and \texttt{-{}-tsv} columns are
(path, name, kind, topics, source, number). The Vim command \texttt{:ImportFact}
(\texttt{<Space>if}) mirrors \texttt{:ImportProblem}.

\subsection*{find-passages and pick-passages (the passage-bank twins)}

The passage bank (section~\ref{sec:ch5-passages}) --- prose shared across
courses, living in a submodule of its own --- is the one bank with
\emph{no metadata block}. The directory is the bank: every \file{.tex} in it is
a passage, and its title and label are derived from its file name rather than
read out of the file. \texttt{-{}-tsv} columns are therefore (path, title), and
what \texttt{-{}-import} emits is a \cs{section} and \cs{label} line, a blank,
then a plain \cs{import}. The Vim command \texttt{:ImportPassage}
(\texttt{<Space>ig}) mirrors the other two.

That one difference accounts for every other. There is nothing to search, so
there are no filters; nothing to inject, so no \texttt{-{}-with}; no membership
to invert, so no \texttt{-{}-untagged}; and no tags, so no \texttt{-{}-vocab}.
All four are refused with a message rather than accepted uselessly.

It is also the one command that \emph{requires} an operand. The bank is a
sibling submodule rather than a directory of the course, so name it ---
\texttt{find-passages PASSAGES/passages}, where \texttt{PASSAGES} stands for
whatever your course checks that submodule out as. \emph{The manual cannot name
it for you}: the checkout directory is the course's own choice, and courses
differ (\file{teaching/} and \file{reusable-teaching-stuff/} are both in use).
Run bare from a course root, with no block to tell a passage from any other
\file{.tex}, it would report every lecture and assessment as a passage --- so it
refuses instead of defaulting to the current directory the way the other four
do.

Under the hood all six commands are thin shims over one shared engine,
\file{find-meta} (\file{find-problems}/\file{find-facts}/\file{find-passages})
and \file{pick-meta}
(\file{pick-problems}/\file{pick-facts}/\file{pick-passages}): a
\texttt{-{}-kind} selects the block
macro (empty for a passage, which has none), the emitted import command and the
key schema; the reading, matching and emitting are written once. A blockless
kind takes its own exit straight after the directory walk, so it never reaches
either of the two block parsers. Run any shim with \texttt{-{}-help} for its concrete
flags.
